\documentclass{article}

\usepackage{float}
\usepackage{amssymb}
\usepackage{amsmath}
\usepackage{hyperref}
\usepackage{graphicx} % Required for inserting images
\usepackage{booktabs} % For professional-looking table rules
\usepackage{caption} % The booktabs package will give you nicer line spacing, but if you don't care about the spacing, you can use the standard hline command
\usepackage{subcaption} % For creating subfigures and subtables
\usepackage{float} % The standard floats.
\usepackage{mathrsfs}
\usepackage{footnote}
\usepackage{algorithm}
\usepackage{algorithmic}
\usepackage{bm}
\usepackage{hyperref}
\hypersetup{
    colorlinks=true,
    linkcolor=red,
    filecolor=red,
    urlcolor=red,
    citecolor=red
    }

\urlstyle{same}
\usepackage[left=4cm, right=4cm, top = 3cm, bottom = 3cm]{geometry}


\title{Phase Diversity Adaptive Optics for Multiphoton Microscopy}
\author{ Max Watzky\footnote{Departments of Physics and Mathematics, Yale University}, \ Shengqi Wang\footnote{Department of Biomedical Engineering, Yale University}, \ and Professor Cristina Rodr\'{\i}guez\footnote{Departments of Applied Physics and Biomedical Engineering, Yale University}}
\date{September 21st, 2026}

\begin{document}
\maketitle

\section{Forward Model}
Let $\mathbf{u}$ be the coordinates of the back focal plane (which we treat as two-dimensional), and
$\mathbf{x}$ be the coordinates of the front focal plane (which we treat as three-dimensional). We denote a 3D stack of images as a map from $\mathbb{R}^3$ to the nonnegative real line:
\begin{equation}
d: \mathbb{R}^3 \to [0, \infty)
\end{equation}
The image stack is generated by convolving a 3D PSF $s_k: \mathbb{R}^3 \to [0, \infty)$ with the underlying 3D object $f: \mathbb{R}^3 \to [0, \infty)$. It is often convenient to consider the individual images separately. We will label the images with the index $k$, and write:
\begin{equation}
d_k(x, y) = \sum_{i=1}^P f(x, y, z_i) \otimes s_k(x, y, z_i - z_k)
\end{equation}
where $z_k$ is the focal plane of the image, and $\{z_i\}$ are the $P$ planes around this focal plane.

The PSF is generated by taking the vector pupil function $\mathbf{H}_k: \mathbb{R}^2 \to \mathbb{C}^3$, applying defocus according to the z-component of $\mathbf{x}$, performing the inverse Fourier transform component by component, and raising the total intensity of the result to the $N$th power:
\begin{align}
\mathbf{h}_k(\mathbf{x}) &= \mathcal{F}^{-1} \{\mathbf{H}_k(\mathbf{u}) e^{i z \gamma(\mathbf{u})}\}, \\
I_k(\mathbf{x}) &= \mathbf{h}_k^\dagger(\mathbf{x})\mathbf{h}_k(\mathbf{x})
= \sum_{c\in\{x,y,z\}} |h_{k,c}(\mathbf{x})|^2, \\
s_k(\mathbf{x}) &= I_k(\mathbf{x})^N.
\end{align}
Here $\mathbf{h}_k=(h_{k,x},h_{k,y},h_{k,z})^{\mathsf T}$ is the electric field, and the dagger denotes the conjugate transpose. The three field components contribute to one intensity image. In particular, we sum their squared magnitudes before raising the result to the $N$th power.
The extra defocus is what sets this model apart from 2D. We calculate it with:
\begin{equation}
\gamma(\mathbf{u}) = \frac{2\pi n_m}{\lambda} \sqrt{1-  \left(\frac{\lambda \|\mathbf{u}\|}{n_m} \right)^2}
\end{equation}
where $n_m$ is the refractive index and $\lambda$ is the vacuum wavelength. In this Fourier representation, $\mathbf{u}$ denotes the lateral spatial frequency associated with a point in the pupil. The same defocus factor is applied to all three field components.
Furthermore, the vector pupil function is generated by a fixed vector pupil mask/apodization $\mathbf{P}(\mathbf{u})$, an underlying phase aberration $\phi(\mathbf{u})$ and a known diversity phase $\theta_k(\mathbf{u})$. In all, we write this as:
\begin{equation}
\mathbf{H}_k(\mathbf{u}) = \mathbf{P}(\mathbf{u}) e^{i \left(\phi(\mathbf{u}) + \theta_k(\mathbf{u})\right)}.
\end{equation}
The underlying phase and the diversity phase are shared by all three components. The vector pupil weights are known and do not depend on the aberration coefficients.

For the fixed, $x$-polarized incident field used here, we can write these weights in terms of the pupil polar angle $\vartheta$ and azimuthal angle $\psi$:
\begin{align}
\sin\vartheta &= \frac{\lambda\|\mathbf{u}\|}{n_m},
\qquad \psi = \arg(u_x+i u_y), \\
\mathbf{P}(\mathbf{u}) &= \frac{C|P(\mathbf{u})|}{\sqrt{\cos\vartheta}}
\begin{pmatrix}
\cos\vartheta+(1-\cos\vartheta)\sin^2\psi \\
(\cos\vartheta-1)\cos\psi\sin\psi \\
-\sin\vartheta\cos\psi
\end{pmatrix}.
\end{align}
Here $P(\mathbf{u})$ is the scalar pupil mask/illumination amplitude, and $C$ includes the fixed propagation prefactor and Fourier-transform normalization. These expressions are used inside the pupil, where $\cos\vartheta>0$; outside the pupil, $\mathbf{P}$ is zero. The angles $\vartheta$ and $\psi$ are distinct from the diversity phase $\theta_k$ and the unknown phase $\phi$. The angular factors account for the redistribution of the incident polarization among the transverse and longitudinal field components.

We assume that the multiphoton excitation depends on the total electric-field intensity as $I_k^N$. We do not include a separate molecular-orientation or polarization-dependent nonlinear response. We also retain the assumptions of a shift-invariant response and incoherent addition of emission from the underlying object, so that the image is a convolution with the real, nonnegative PSF $s_k$.
Sometimes, we will expand $\phi(\mathbf{u})$ in a basis $\{ Z_i\}$. In this case, we will use the vector $\bm{\alpha}$ to describe the coefficients of $\phi$ in this basis.

\section{Cost Functionals}
Let $\{d_k\}$ be the set of acquired images, with known diversity phases $\{\theta_k\}$. The goal is to estimate the underlying aberration $\phi(\mathbf{u})$ and the underlying sample $f(\mathbf{x})$ to minimize a cost functional $\mathcal{L} \{f(\mathbf{x}), \bm{\alpha}\}$. In keeping with the literature, we will treat two cases: Gaussian and Poisson noise. In the Gaussian case, we seek to minimize the least-squares objective:
\begin{equation}
\mathcal{L}_g \{f, \bm{\alpha}\} = \sum_{k=1}^K \sum_{(x,y)} \left[d_k(x, y) - \mu_k(x,y) \right]^2
\label{eq:gaussian_cost_functional}
\end{equation}
where we use the shorthand:
\begin{align}
\mu_k(x, y) = \sum_{i=1}^P f(x, y, z_i) \otimes s_k(x, y, z_i - z_k)
\end{align}
where $s_k$ here denotes the prediction of the PSF from $\bm{\alpha}$. We can also write out the Poisson objective:
\begin{equation}
\mathcal{L}_p \{f, \bm{\alpha}\} = 2 \sum_{k=1}^K \sum_{(x,y)} \left[d_k(x, y) \ln(d_k(x,y)) - d_k(x,y) \ln(\mu_k(x,y)) - (d_k(x,y) - \mu_k(x,y))\right]
\end{equation}

\section{Gaussian Noise}
We will first attempt to solve the problem for the case of Gaussian noise.

\subsection{Object Estimation}
We begin by estimating the underlying object $f$. By Parseval's theorem, we can rewrite the Gaussian cost functional in \autoref{eq:gaussian_cost_functional} in the Fourier domain:
\begin{equation}
\mathcal{L}_g \{F, \bm{\alpha}\} = \sum_{k=1}^K \sum_{(u_x, u_y)} \left| D_k - \sum_{i=1}^P F_{z_i} S_k (z_i - z_k)\right|^2
\label{eq:gaussian_cost_functional_fourier}
\end{equation}
where $D_k$ denotes the Fourier transform of $d_k$, $F_{z_i}$ denotes the Fourier transform of $f(x, y, z_i)$ and $S_k(z_i - z_k)$ denotes the Fourier transform of the $k$th simulated PSF slice at $z_i - z_k$. The notation $\mathcal{F}_{xy}$ used below denotes the same lateral Fourier transform.
At each fixed spatial frequency, the unconstrained object that minimizes the cost functional will satisfy $\partial \mathcal{L}_g/ \partial F_{z_j}^* = 0$ for each $j=1,\ldots,P$. We compute each of these w.r.t the complex conjugate of $F_{z_j}$ so that the resulting expression will include $F_{z_j}$ directly:
\begin{align}
\frac{\partial \mathcal{L}_g}{\partial F_{z_j}^*} &= \frac{\partial}{\partial F_{z_j}^*} \sum_{k=1}^K \sum_{(u_x, u_y)} \left| D_k - \sum_{i=1}^P F_{z_i} S_k (z_i - z_k)\right|^2 \\
&=  \sum_{k=1}^K \frac{\partial}{\partial F_{z_j}^*} \left[ \left(D_k - \sum_{i=1}^P F_{z_i} S_k (z_i - z_k)\right) \left(D^*_k - \sum_{i=1}^P F^*_{z_i} S^*_k (z_i - z_k)\right) \right] \\
&=  - \sum_{k=1}^K \left[ S_k^*(z_j- z_k) \left(D_k - \sum_{i=1}^P F_{z_i} S_k (z_i - z_k)\right) \right]
\end{align}
Rearranging, we get:
\begin{align}
\sum_{i=1}^P \left[\sum_{k=1}^K S_k^*(z_j - z_k) S_k(z_i - z_k)  \right] F_{z_i} = \sum_{k=1}^K S_k^*(z_j - z_k) D_k
\label{eq:optimize_for_f}
\end{align}
We get $P$ equations of this form, one for each index $j=1,...,P$. Together, they form the matrix equation:
\begin{align}
(\mathbf{S}^\dagger \mathbf{S}) \mathbf{F} = \mathbf{S}^\dagger \mathbf{D}
\end{align}
It is helpful to briefly remind ourselves of the dimensions of these matrices: $\mathbf{F}$ is a single column with one row for each equation ($P$). By similar logic, we can infer that $\mathbf{D}$ has one column and one row for each image ($K$). $\mathbf{S}$ is $K \times P$, and $\mathbf{S}^\dagger$ is $P \times K$, so, $\mathbf{S}^\dagger \mathbf{S}$ is $P \times P$.
We can thus verify the dimensions of our equation are correct. Solving, we get the estimator:
\begin{align}
\hat{\mathbf{F}} = (\mathbf{S}^\dagger \mathbf{S})^{-1} \mathbf{S}^\dagger \mathbf{D}
\end{align}
where we have assumed that $\mathbf{S}^\dagger \mathbf{S}$ is invertible. This means that all of the columns of $\mathbf{S}$ must be linearly independent. Since $\mathbf{S}$ is $K \times P$, it must be the case that $K \ge P$; otherwise this is impossible. As Kner points out, this tells us that we must have at least $P$ measurements to estimate the image at $P$ planes.

Sometimes, we would like an estimate, even when we have insufficient data. We therefore introduce a fixed regularization parameter $\rho>0$ as follows:
\begin{align}
\hat{\mathbf{F}} = (\mathbf{S}^\dagger \mathbf{S} + \rho I_{P})^{-1} \mathbf{S}^\dagger \mathbf{D}
\label{eq:regularized_F}
\end{align}
which will allow us to solve for the inverse even if $\mathbf{S}$ has insufficient rank. This is equivalent to placing a squared $L^2$ norm penalty on $\mathbf{F}$ in the original cost function \autoref{eq:gaussian_cost_functional_fourier}.

\subsection{Aberration Estimation}
Next, we will discuss the process of estimating the aberration along a basis $\{Z_i\}$, with real coefficients $\alpha_i$:
\begin{equation}
\phi(\mathbf{u}) = \sum_i \alpha_i Z_i(\mathbf{u})
\end{equation}
Here $\alpha_i$ multiplies the dimensionless basis function directly in the phase. If a coefficient is instead stored in waves, its phase coefficient is $2\pi$ times its wave strength; derivatives with respect to wave strengths therefore use $2\pi Z_i$ in place of $Z_i$.
We will begin by defining the following variable, the residual between our $k$th estimated image (which we generate from the estimated object $\hat{\mathbf{F}}$) and the actual image $D_k$:
\begin{equation}
r_k = D_k - \sum_{i=1}^P \hat{F}_{z_i} S_{k}(z_i - z_k)
\end{equation}
With $\hat{\mathbf{F}}$ plugged in, and the additional squared $L^2$ penalty added, we rewrite the original cost functional \autoref{eq:gaussian_cost_functional_fourier} as:
\begin{equation}
\mathcal{L}_g \{\bm{\alpha}\} = \sum_{(u_x, u_y)} \sum_{k=1}^K r_k r_k^*  + \rho \sum_{(u_x, u_y)} \sum_{i=1}^P \hat{F}_{z_i} \hat{F}^*_{z_i}
\end{equation}
To take the gradient $\nabla \mathcal{L}_g$, we begin by taking the derivative with respect to an arbitrary aberration coefficient $\alpha_m$:
\begin{align}
\frac{\partial \mathcal{L}_g}{\partial \alpha_m} &= \sum_{(u_x, u_y)} \sum_{k=1}^K \left[r_k^* \left(\frac{\partial r_k}{\partial \alpha_m} \right)  + r_k \left(\frac{\partial r_k}{\partial \alpha_m} \right)^*\right] \nonumber\\
&\quad + \rho \sum_{(u_x, u_y)} \sum_{i=1}^P \left[\hat{F}_{z_i}^* \left(\frac{\partial \hat{F}_{z_i}}{\partial \alpha_m} \right)  + \hat{F}_{z_i} \left(\frac{\partial \hat{F}_{z_i}}{\partial \alpha_m} \right)^*\right]
\end{align}
In both terms, we recognize sums of the form $w + w^* = 2 \Re(w)$. Using this identity, we rewrite:
\begin{equation}
\frac{\partial \mathcal{L}_g}{\partial \alpha_m} = 2 \Re \left[\sum_{(u_x, u_y)} \sum_{k=1}^K r_k^* \frac{\partial r_k}{\partial \alpha_m} \right] + 2 \rho \Re \left[\sum_{(u_x, u_y)} \sum_{i=1}^P \hat{F}_{z_i}^* \frac{\partial \hat{F}_{z_i}}{\partial \alpha_m} \right]
\end{equation}
We examine the next derivative:
\begin{equation}
\frac{\partial r_k}{\partial \alpha_m} = \frac{\partial}{\partial \alpha_m} \left[D_k - \sum_{i=1}^P \hat{F}_{z_i} S_k(z_i - z_k) \right] = - \sum_{i=1}^P \left[\frac{\partial \hat{F}_{z_i}}{\partial \alpha_m} S_k(z_i - z_k) + \hat{F}_{z_i} \frac{\partial S_k(z_i - z_k)}{\partial \alpha_m} \right]
\end{equation}
From here on, for compactness we denote:
\begin{align}
S_{ki} = S_k(z_i - z_k)
\end{align}
Plugging in our results, we get:
\begin{align}
\frac{\partial \mathcal{L}_g}{\partial \alpha_m} &= -2 \Re \left[\sum_{(u_x, u_y), k, i} r_k^* S_{ki} \frac{\partial \hat{F}_{z_i}}{\partial \alpha_m}\right] \nonumber\\
&\quad - 2 \Re \left[ \sum_{(u_x, u_y), k, i} r_k^* \hat{F}_{z_i} \frac{\partial S_{ki}}{\partial \alpha_m} \right] \nonumber\\
&\quad + 2 \rho \Re \left[\sum_{(u_x, u_y), i} \hat{F}^*_{z_i} \frac{\partial \hat{F}_{z_i}}{\partial \alpha_m}\right]
\end{align}
We briefly examine the first and third terms, since they both involve the derivative of the optimized object with respect to the aberration coefficient. Inside the real-part operator, we have:
\begin{align}
-\sum_{(u_x, u_y)} \sum_i \left[\sum_k r_k^* S_{ki} \right] \frac{\partial \hat{F}_{z_i}}{\partial \alpha_m} + \rho \sum_{(u_x, u_y)} \sum_i \hat{F}^*_{z_i} \frac{\partial \hat{F}_{z_i}}{\partial \alpha_m}
\end{align}
However, recall from \autoref{eq:regularized_F} that:
\begin{align}
(\mathbf{S}^\dagger \mathbf{S} + \rho I_P) \hat{\mathbf{F}} = \mathbf{S}^\dagger (\mathbf{r} + \mathbf{S} \hat{\mathbf{F}})
\end{align}
which gives:
\begin{align}
\rho \hat{\mathbf{F}} = \mathbf{S}^\dagger \mathbf{r}
\end{align}
which allows us to cancel the corresponding terms in the summation. The result is:
\begin{align}
\frac{\partial \mathcal{L}_g}{\partial \alpha_m} = - 2 \Re \left[ \sum_{(u_x, u_y), k, i} r_k^* \hat{F}_{z_i} \frac{\partial S_{ki}}{\partial \alpha_m} \right]
\end{align}
The next task is to evaluate the PSF derivative. Recall the forward model:
\begin{align}
\frac{\partial S_{ki}}{\partial \alpha_m} = \frac{\partial}{\partial \alpha_m} \mathcal{F} \{s_k(x, y, z_i - z_k) \} = \mathcal{F} \left\{ \frac{\partial s_k(x, y, z_i - z_k)}{\partial \alpha_m} \right\}
\end{align}
where, writing $I_k=\mathbf{h}_k^\dagger\mathbf{h}_k$, we have:
\begin{align}
\frac{\partial I_k}{\partial \alpha_m}
&= \frac{\partial}{\partial \alpha_m}
\sum_{c\in\{x,y,z\}} h_{k,c}h_{k,c}^* \\
&= 2\Re\left[\sum_{c\in\{x,y,z\}}
h_{k,c}^*\frac{\partial h_{k,c}}{\partial\alpha_m}\right]
=2\Re\left[\mathbf{h}_k^\dagger
\frac{\partial\mathbf{h}_k}{\partial\alpha_m}\right], \\
\frac{\partial s_k(x,y,z_i-z_k)}{\partial\alpha_m}
&= N I_k(z_i-z_k)^{N-1}
\frac{\partial I_k(z_i-z_k)}{\partial\alpha_m} \\
&= 2N I_k(z_i-z_k)^{N-1}
\Re\left[\mathbf{h}_k^\dagger(z_i-z_k)
\frac{\partial\mathbf{h}_k(z_i-z_k)}{\partial\alpha_m}\right].
\end{align}
The field derivative is evaluated component by component:
\begin{align}
\frac{\partial\mathbf{h}_k}{\partial\alpha_m}
&= \frac{\partial}{\partial\alpha_m}
\mathcal F^{-1}\{\mathbf H_k(\mathbf u)e^{iz\gamma(\mathbf u)}\} \\
&= \mathcal F^{-1}\left\{iZ_m(\mathbf u)\mathbf H_k(\mathbf u)e^{iz\gamma(\mathbf u)}\right\},
\end{align}
where we have used that the only $\alpha_m$ dependence in this term is in $\mathbf H_k$, where:
\begin{align}
\frac{\partial\mathbf H_k}{\partial\alpha_m}
&= \frac{\partial}{\partial\alpha_m}
\mathbf P(\mathbf u)e^{i(\phi(\mathbf u)+\theta_k(\mathbf u))}
=iZ_m\mathbf H_k
\end{align}
by expanding $\phi$ in our basis as $\sum_j Z_j \alpha_j$. Summarizing, we have:
\begin{equation}
\begin{aligned}
\frac{\partial \mathcal{L}_g}{\partial \alpha_m} ={}& -4N\,\Re \sum_{(u_x,u_y)} \sum_{k=1}^{K} \left[D_k-\sum_{i=1}^{P}\hat{F}_{z_i}S_k(z_i-z_k)\right]^*
\\&\quad\times\left[\sum_{j=1}^{P}\hat{F}_{z_j}\,\mathcal{F}_{xy}\left\{I_k(z_j-z_k)^{N-1}
\Re\left[\mathbf h_k^\dagger(z_j-z_k)\,
\mathcal{F}^{-1}\left\{iZ_m\mathbf H_k e^{i(z_j-z_k)\gamma}\right\}\right]\right\}\right].
\end{aligned}
\end{equation}
To compute the Hessian, we also need the second derivatives of the vector field and intensity. The second field derivative is:
\begin{equation}
\frac{\partial^2\mathbf h_k}{\partial\alpha_m\partial\alpha_\ell}
=\mathcal F^{-1}\left\{-Z_mZ_\ell\mathbf H_k e^{iz\gamma}\right\}.
\end{equation}
Differentiating the intensity once more, we get:
\begin{equation}
\begin{aligned}
\frac{\partial^2 I_k}{\partial\alpha_m\partial\alpha_\ell}
=2\Re\Bigg[&
\left(\frac{\partial\mathbf h_k}{\partial\alpha_m}\right)^\dagger
\frac{\partial\mathbf h_k}{\partial\alpha_\ell} \\
&+\mathbf h_k^\dagger
\frac{\partial^2\mathbf h_k}{\partial\alpha_m\partial\alpha_\ell}
\Bigg].
\end{aligned}
\end{equation}
Applying the chain rule to $s_k=I_k^N$, we obtain:
\begin{equation}
\begin{aligned}
\frac{\partial^2s_k}{\partial\alpha_m\partial\alpha_\ell}
={}&N(N-1)I_k^{N-2}
\frac{\partial I_k}{\partial\alpha_m}
\frac{\partial I_k}{\partial\alpha_\ell} \\
&+NI_k^{N-1}
\frac{\partial^2I_k}{\partial\alpha_m\partial\alpha_\ell}.
\end{aligned}
\end{equation}
The terms proportional to $N-1$ are omitted when $N=1$. The component sums are taken before forming the product of intensity derivatives, just as they are taken before raising the intensity to the $N$th power in the forward model.

The Hessian of the cost functional, including the dependence of the estimated object on the aberration, takes the following form:
\begin{equation}
\begin{aligned}
\frac{\partial^2 \mathcal{L}_g}
{\partial\alpha_m\,\partial\alpha_\ell}
={}&
8N^2\operatorname{Re}
\sum_{(u_x,u_y)}\sum_{k=1}^{K}
\left[
\sum_{i=1}^{P}\widehat F_{z_i}\,
\mathcal F_{xy}
\left\{
I_k(z_i-z_k)^{N-1}a_{ki,m}
\right\}
\right]^*
\\
&\quad\times
\left[
\sum_{j=1}^{P}\widehat F_{z_j}\,
\mathcal F_{xy}
\left\{
I_k(z_j-z_k)^{N-1}a_{kj,\ell}
\right\}
\right]
\\[4pt]
&-
8N(N-1)\operatorname{Re}
\sum_{(u_x,u_y)}\sum_{k=1}^{K}
r_k^*
\sum_{j=1}^{P}\widehat F_{z_j}\,
\mathcal F_{xy}
\left\{
I_k(z_j-z_k)^{N-2}a_{kj,m}a_{kj,\ell}
\right\}
\\[4pt]
&-
4N\operatorname{Re}
\sum_{(u_x,u_y)}\sum_{k=1}^{K}
r_k^*
\sum_{j=1}^{P}\widehat F_{z_j}\,
\mathcal F_{xy}
\left\{
I_k(z_j-z_k)^{N-1}b_{kj,m\ell}
\right\}
\\[4pt]
&-
2\operatorname{Re}
\sum_{(u_x,u_y)}
\sum_{i=1}^{P}\sum_{j=1}^{P}
c_{i,m}^*
\left[
\left(\mathbf S^\dagger\mathbf S+\rho I_P\right)^{-1}
\right]_{ij}
c_{j,\ell}.
\end{aligned}
\end{equation}
where:
\begin{align*}
a_{ki,m}
&=
\operatorname{Re}\left[
\mathbf h_k^\dagger(z_i-z_k)\,
\mathcal F^{-1}
\left\{
iZ_m\mathbf H_k e^{i(z_i-z_k)\gamma}
\right\}
\right],
\\[6pt]
b_{ki,m\ell}
&=
\operatorname{Re}\Bigg[
\left(
\mathcal F^{-1}
\left\{
Z_\ell\mathbf H_k e^{i(z_i-z_k)\gamma}
\right\}
\right)^\dagger
\mathcal F^{-1}
\left\{
Z_m\mathbf H_k e^{i(z_i-z_k)\gamma}
\right\}
\\
&\qquad\qquad-
\mathbf h_k^\dagger(z_i-z_k)\,
\mathcal F^{-1}
\left\{
Z_mZ_\ell\mathbf H_k e^{i(z_i-z_k)\gamma}
\right\}
\Bigg],
\\[6pt]
c_{j,m}
&=
2N\sum_{k=1}^{K}\Bigg[
\left(
\mathcal F_{xy}
\left\{
I_k(z_j-z_k)^{N-1}a_{kj,m}
\right\}
\right)^*r_k
\\
&\qquad\qquad-
S_k^*(z_j-z_k)
\sum_{i=1}^{P}\widehat F_{z_i}\,
\mathcal F_{xy}
\left\{
I_k(z_i-z_k)^{N-1}a_{ki,m}
\right\}
\Bigg],
\\[6pt]
r_k
&=
D_k-\sum_{i=1}^{P}\widehat F_{z_i}S_k(z_i-z_k),
\\[6pt]
\mathbf h_k(z)
&=
\mathcal F^{-1}\left\{\mathbf H_k e^{iz\gamma}\right\},
\\[6pt]
I_k(z)
&=
\mathbf h_k^\dagger(z)\mathbf h_k(z)
=\sum_{c\in\{x,y,z\}}|h_{k,c}(z)|^2,
\\[6pt]
S_k(z)
&=
\mathcal F_{xy}\left\{I_k(z)^N\right\},
\\[6pt]
[\mathbf S]_{ki}
&=
S_k(z_i-z_k),
\\[6pt]
\mathbf D
&=
(D_1,\ldots,D_K)^{\mathsf T},
\\[6pt]
\widehat{\mathbf F}
&=
(\widehat F_{z_1},\ldots,\widehat F_{z_P})^{\mathsf T}
=
\left(\mathbf S^\dagger\mathbf S+\rho I_P\right)^{-1}
\mathbf S^\dagger\mathbf D,
\\[6pt]
\mathcal{L}_g(\bm{\alpha})
&=
\sum_{(u_x,u_y)}
\left[
\sum_{k=1}^{K}|r_k|^2
+\rho\sum_{i=1}^{P}|\widehat F_{z_i}|^2
\right].
\end{align*}
The inner products in $a_{ki,m}$ and $b_{ki,m\ell}$ sum over the three field components. The quantities $S_k$, $r_k$ and $c_{j,m}$ remain scalars at each lateral spatial frequency. Thus $\mathbf S$ remains a $K\times P$ matrix, with one row per acquired intensity image.

The last Hessian term accounts for the change in the estimated object. To see this, differentiate the regularized object equation at each fixed spatial frequency:
\begin{equation}
\left(\mathbf S^\dagger\mathbf S+\rho I_P\right)
\frac{\partial\widehat{\mathbf F}}{\partial\alpha_m}
=
\left(\frac{\partial\mathbf S}{\partial\alpha_m}\right)^\dagger\mathbf r
-\mathbf S^\dagger\frac{\partial\mathbf S}{\partial\alpha_m}\widehat{\mathbf F}
=\mathbf c_m,
\end{equation}
where $\mathbf c_m=(c_{1,m},\ldots,c_{P,m})^{\mathsf T}$. This gives the same object-response term as above, with the PSFs and their derivatives computed from the vector fields.

Finally, writing $\mathbf g=\nabla_{\bm\alpha}\mathcal L_g$ and $\mathbf H=\nabla_{\bm\alpha}^2\mathcal L_g$, the Newton step is:
\begin{align}
\Delta \mathbf{\alpha} = -\mathbf{H}^{-1} \mathbf{g}
\end{align}
\end{document}
